%------------------------------------------------------------------------------
% Chapter 3: Methodology
%------------------------------------------------------------------------------

\chapter{Methodology}\label{chap:methodology}

This chapter describes the implementation used to study automated algorithm
design under stochastic fitness evaluation. The method builds on the Large
Language Model Evolutionary Algorithm (LLaMEA), in which a language model
generates executable optimization algorithms and iteratively modifies them
using information obtained from their evaluation \cite{LLaMEA2024}. In this
thesis, the LLaMEA process is combined with a controlled stochastic
evaluation environment and locally executed large language models (LLMs).

The central methodological principle is to separate the optimization
environment experienced by a generated algorithm from the information used
by the synthesis framework to assess that algorithm. An LLM receives the
optimization task, the current synthesis context, and the selected prompt
configuration. It generates a Python implementation of an optimization
heuristic. The generated algorithm is then extracted, checked, and executed
in an isolated environment. During this execution, the optimizer interacts
with the objective only through the provided evaluation interface. Under
stochastic conditions, this interface returns noisy objective values, while
the clean objective remains available inside the evaluation framework for
diagnostic assessment.

The synthesis process is therefore not a direct search over numerical
solutions. Instead, the evolutionary search operates over executable
optimization algorithms. The language model acts as the semantic variation
operator, while the surrounding framework is responsible for generating,
executing, evaluating, and selecting these algorithms. The information
obtained from each candidate is subsequently used to update the synthesis
state and construct the context for later generations.

The complete process can be summarized as

\begin{equation}
\begin{aligned}
\text{Optimization task}
&\rightarrow
\text{Synthesis context}
\rightarrow
\text{LLM generation}
\\
&\rightarrow
\text{Code extraction}
\rightarrow
\text{Validation}
\rightarrow
\text{Isolated execution}
\\
&\rightarrow
\text{Objective evaluation}
\rightarrow
\text{Clean diagnostic evaluation}
\\
&\rightarrow
\text{Feedback}
\rightarrow
\text{Parent comparison}
\rightarrow
\text{Updated incumbent}.
\end{aligned}
\label{eq:method_overview}
\end{equation}

Figure~\ref{fig:llamea_synthesis_overview} provides a corresponding overview
of the synthesis loop. Starting from the optimization task and synthesis
configuration, the LLM generates a candidate optimization algorithm. The
candidate is then extracted, validated, and executed by the evaluation
framework. Its performance is evaluated through the objective interface,
while the returned solution is additionally assessed using the clean
objective for diagnostic purposes. The resulting information is used for
parent--child comparison and for constructing the synthesis context of the
next generation. The process continues until the synthesis run terminates.

\begin{figure}[htbp]
\centering
\resizebox{\textwidth}{!}{%
\begin{tikzpicture}[
    font=\footnotesize,
    >=Stealth,
    ioNode/.style={
        rectangle,
        draw=black!80,
        fill=yellow!15,
        text width=13cm,
        align=center,
        rounded corners=4pt,
        inner sep=8pt,
        thick
    },
    processNode/.style={
        rectangle,
        draw=black!80,
        fill=blue!5,
        text width=5.7cm,
        align=center,
        rounded corners=4pt,
        inner sep=7pt,
        thick,
        minimum height=1.45cm
    },
    evalNode/.style={
        processNode,
        fill=blue!10
    },
    selectionNode/.style={
        rectangle,
        draw=black!80,
        fill=purple!10,
        text width=12.5cm,
        align=center,
        rounded corners=4pt,
        inner sep=8pt,
        thick,
        minimum height=1.6cm
    },
    decision/.style={
        diamond,
        draw=black!80,
        fill=gray!15,
        thick,
        aspect=3.0,
        inner sep=4pt,
        align=center,
        font=\scriptsize,
        minimum width=5.8cm
    },
    outNode/.style={
        rectangle,
        draw=black!80,
        fill=green!15,
        text width=13cm,
        align=center,
        rounded corners=4pt,
        inner sep=8pt,
        thick,
        font=\footnotesize\bfseries
    },
    arrow/.style={
        ->,
        >=Stealth,
        thick,
        draw=black!75
    },
    fbArrow/.style={
        ->,
        >=Stealth,
        thick,
        draw=purple!80!black
    },
    elbl/.style={
        font=\scriptsize,
        fill=white,
        inner sep=2pt,
        rounded corners=1pt
    },
    elblB/.style={
        elbl,
        font=\scriptsize\bfseries
    }
]

%----------------------------------------------------------------------
% Input
%----------------------------------------------------------------------

\node[ioNode] (inputs) {
    \textbf{Optimization Task and Synthesis Configuration} \\[3pt]
    \scriptsize
    Continuous black-box problem
    \quad $\bullet$ \quad
    Prompt configuration
    \quad $\bullet$ \quad
    Synthesis budget
};

%----------------------------------------------------------------------
% Synthesis context
%----------------------------------------------------------------------

\node[
    processNode,
    below=0.8cm of inputs
] (context) {
    \textbf{Synthesis Context} \\[3pt]
    \scriptsize
    Optimization task, previous candidate information,
    incumbent, and prompt configuration
};

%----------------------------------------------------------------------
% Generation
%----------------------------------------------------------------------

\node[
    processNode,
    below=0.8cm of context
] (generation) {
    \textbf{LLM Candidate Generation} \\[3pt]
    \scriptsize
    Generates executable optimization heuristic
    $\mathcal{A}_{\text{child}}$
};

%----------------------------------------------------------------------
% Code extraction and validation
%----------------------------------------------------------------------

\node[
    processNode,
    below=0.8cm of generation
] (validation) {
    \textbf{Code Extraction and Validation} \\[3pt]
    \scriptsize
    Extracts generated source code and checks
    syntactic and structural constraints
};

%----------------------------------------------------------------------
% Execution
%----------------------------------------------------------------------

\node[
    processNode,
    below=0.8cm of validation
] (execution) {
    \textbf{Isolated Candidate Execution} \\[3pt]
    \scriptsize
    Runs the candidate under a controlled execution
    environment and timeout
};

%----------------------------------------------------------------------
% Evaluation: two parallel components
%----------------------------------------------------------------------

\node[
    evalNode,
    below left=1.0cm and 0cm of execution.south,
    anchor=north west
] (stochastic) {
    \textbf{Stochastic Objective Evaluation} \\[3pt]
    \scriptsize
    Candidate interacts with the objective through
    noisy fitness observations
};

\node[
    evalNode,
    right=1.0cm of stochastic
] (clean) {
    \textbf{Clean Diagnostic Evaluation} \\[3pt]
    \scriptsize
    Returned solution is assessed using the
    underlying clean objective
};

%----------------------------------------------------------------------
% Selection and feedback
%----------------------------------------------------------------------

\node[
    selectionNode,
    below=1.0cm of $(stochastic.south)!0.5!(clean.south)$
] (selection) {
    \textbf{Parent--Child Comparison and Context Update} \\[3pt]
    \scriptsize
    Compares candidate performance with the incumbent,
    updates the retained solution, and constructs feedback
    for the next generation
};

%----------------------------------------------------------------------
% Termination
%----------------------------------------------------------------------

\node[
    decision,
    below=1.0cm of selection
] (termination) {
    \textbf{Synthesis Terminated?}
};

%----------------------------------------------------------------------
% Output
%----------------------------------------------------------------------

\node[
    outNode,
    below=1.0cm of termination
] (output) {
    \textbf{Final Output: Synthesized Optimization Heuristic}
    \\[3pt]
    \scriptsize
    Retained incumbent algorithm and associated synthesis information
};

%----------------------------------------------------------------------
% Loop boundary
%----------------------------------------------------------------------

\begin{scope}[on background layer]
    \node[
        draw=gray!50,
        dashed,
        fill=gray!3,
        rounded corners=8pt,
        inner sep=12pt,
        fit=(context) (generation) (validation) (execution)
             (stochastic) (clean) (selection),
        label={[
            font=\bfseries\scriptsize,
            gray!70!black
        ]above:Iterative LLaMEA Synthesis Loop}
    ] (loopGroup) {};
\end{scope}

%----------------------------------------------------------------------
% Forward arrows
%----------------------------------------------------------------------

\draw[arrow]
    (inputs.south)
    -- (context.north);

\draw[arrow]
    (context.south)
    -- (generation.north);

\draw[arrow]
    (generation.south)
    -- node[elblB, right=2pt]
    {Candidate program}
    (validation.north);

\draw[arrow]
    (validation.south)
    -- (execution.north);

%----------------------------------------------------------------------
% Evaluation arrows
%----------------------------------------------------------------------

\draw[arrow]
    (execution.south)
    -- ++(0,-0.35)
    -| (stochastic.north);

\draw[arrow]
    (execution.south)
    -- ++(0,-0.35)
    -| (clean.north);

\draw[arrow]
    (stochastic.south)
    -- ++(0,-0.45)
    -| (selection.west);

\draw[arrow]
    (clean.south)
    -- ++(0,-0.45)
    -| (selection.east);

%----------------------------------------------------------------------
% Termination
%----------------------------------------------------------------------

\draw[arrow]
    (selection.south)
    -- (termination.north);

\draw[arrow]
    (termination.south)
    -- node[elblB, right=2pt]
    {Yes}
    (output.north);

%----------------------------------------------------------------------
% Next-generation feedback
%----------------------------------------------------------------------

\draw[fbArrow]
    (termination.west)
    -- ++(-2.8,0)
    node[
        elblB,
        pos=0.5,
        above,
        text=purple!80!black
    ] {No: next generation}
    |- ($(context.west)+(-0.6,0)$);

\end{tikzpicture}
}
\caption{Overview of the LLaMEA-based algorithm synthesis process used in
this thesis. The model generates candidate optimization heuristics, which
are evaluated and used to update the synthesis context for subsequent
generations.}
\label{fig:llamea_synthesis_overview}
\end{figure}

%==============================================================================
\section{LLM-Based Algorithm Synthesis}
\label{sec:llm_algorithm_synthesis}
%==============================================================================

Each synthesis run starts with a continuous black-box optimization task and a
prompt template defining the programming interface available to the generated
algorithm. The generated algorithm is represented as a Python class following
the interface used by LLaMEA. Rather than producing a numerical candidate
solution directly, the model therefore produces a complete optimization
procedure. This procedure may contain its own initialization strategy,
population management, search operators, parameter updates, and stopping
logic.

The object being evolved is consequently not a solution vector but an
executable optimization algorithm. The LLM can modify the structure of the
search procedure itself, allowing changes to affect how candidate solutions
are generated, how information from previous evaluations is used, and how
the search is adapted throughout the optimization run.

This differs from conventional evolutionary optimization, where variation is
usually applied to numerical parameters or solution vectors. In the present
setting, the language model performs the variation at the program level. The
surrounding synthesis framework then determines whether the generated
program is valid and whether the resulting optimization behavior is
sufficiently effective for the candidate to be retained.


%------------------------------------------------------------------------------
\subsection{Synthesis Context and Candidate Generation}
\label{subsec:synthesis_context}
%------------------------------------------------------------------------------

The synthesis context contains the information required by the model to
generate or modify an algorithm. In the elitist $(1+1)$ setting, the context
can include the task description, information about previously generated
algorithms, the current incumbent, and feedback obtained from earlier
evaluations \cite{LLaMEA2024}. The context can be represented conceptually
as

\begin{equation}
\mathcal{C}_g =
\left\{
\mathcal{P}_{\mathrm{task}},
\mathcal{H}_{1:g-1},
\mathcal{A}_{\mathrm{inc}},
\mathcal{F}_{\mathrm{prompt}}
\right\},
\label{eq:synthesis_context}
\end{equation}

where $\mathcal{P}_{\mathrm{task}}$ describes the optimization problem,
$\mathcal{H}_{1:g-1}$ contains relevant information from previous
generations, $\mathcal{A}_{\mathrm{inc}}$ denotes the current incumbent
algorithm, and $\mathcal{F}_{\mathrm{prompt}}$ represents the additional
instructions defined by the selected prompt configuration.

The synthesis history can contain information about previous candidate
algorithms, their observed performance, and execution failures. This
information gives the model a basis for modifying previously generated
approaches rather than generating every candidate independently. The
incumbent provides the current reference algorithm, while the previous
evaluation history provides additional information about approaches that have
already been explored.

The LLM uses this context to generate a new candidate. The model can modify
an existing algorithm or produce a substantially different search procedure.
Changes can therefore affect the complete search strategy rather than only
individual parameter values. For example, a candidate can introduce a
different population update mechanism, alter the balance between exploration
and exploitation, adapt its parameters during the run, or change the way
candidate solutions are generated.

The generated response may contain explanatory text in addition to Python
source code. The implementation consequently extracts the candidate source
before execution. The resulting program is then passed to the validation
stage rather than being executed directly from the raw model response.


%------------------------------------------------------------------------------
\subsection{Code Extraction and Validation}
\label{subsec:code_validation}
%------------------------------------------------------------------------------

The code extraction stage identifies the Python source contained in the
generated response and prepares it for execution. Since language model
responses can contain explanations, Markdown formatting, or other text in
addition to the requested program, the extraction step provides a clear
boundary between the generation process and the executable candidate.

The resulting program is checked structurally using Abstract Syntax Tree
(AST) analysis. This validation stage detects malformed or disallowed code
before it reaches the execution environment. The implementation permits the
numerical functionality required for the optimization task while preventing
generated code from delegating the search to external optimization routines
or interacting directly with the surrounding operating system.

The purpose of this stage is not to constrain the model to one particular
optimization strategy. Instead, it ensures that different generated
algorithms can be evaluated under the same interface and within the same
execution constraints. Candidates that fail the required structural checks
are rejected before they can affect the synthesis state.

After validation, the candidate is compiled and executed in an isolated
worker process. Isolation prevents variables or state created by one
candidate from affecting the surrounding synthesis process. It also allows
the execution of a generated algorithm to be terminated independently if it
does not complete within the permitted execution time.

The execution environment additionally checks the numerical validity of the
returned solution. Invalid returned solutions, including non-finite
coordinates or coordinates outside the allowed search domain, are rejected
rather than being treated as successful optimization results.

Runtime failures are handled within the same evaluation stage. A generated
algorithm may fail because of numerical operations, incorrect array
dimensions, invalid use of the objective interface, or other implementation
errors. Such failures are recorded rather than terminating the complete
synthesis process. Where appropriate, the corresponding error information
can be included in the context supplied to later generations. This allows
the model to use execution feedback when modifying a previously generated
program, consistent with the feedback-driven design of LLaMEA
\cite{LLaMEA2024}.


%------------------------------------------------------------------------------
\subsection{Isolated Candidate Execution}
\label{subsec:isolated_execution}
%------------------------------------------------------------------------------

After passing validation, the candidate is executed independently from the
synthesis controller. The generated optimizer interacts with the optimization
problem through the objective-function interface provided to it. The
optimizer does not directly access the analytical benchmark function or the
known optimum.

This maintains the black-box setting of the optimization problem. From the
perspective of the generated algorithm, the objective is an externally
provided function that can be queried at candidate points. The optimizer
must therefore construct its search strategy from the fitness values
returned by these queries.

The isolated worker process also provides a boundary between candidate
execution and the rest of the synthesis framework. A candidate that enters
an infinite loop, exceeds the permitted runtime, or encounters a runtime
failure can be terminated without stopping the complete synthesis process.
The execution result can consequently be classified as either a valid
optimization result or a failure with corresponding diagnostic information.

For a successful candidate, the framework obtains the solution returned by
the optimizer together with the information required to assess its quality.
For an unsuccessful candidate, the framework records the type of failure and
the available diagnostic information. In both cases, the result becomes
part of the evolutionary feedback process.

This execution procedure also ensures that all candidates are evaluated
under the same basic computational constraints. Differences in performance
therefore arise from the generated optimization procedure and its interaction
with the evaluation environment rather than from different execution
interfaces.


%==============================================================================
\section{Stochastic Fitness Evaluation}
\label{sec:stochastic_evaluation}
%==============================================================================

The main adaptation introduced in this thesis is the stochastic objective
evaluation used during candidate execution. The generated optimizer does not
receive direct access to the analytical form of the benchmark function or to
its optimum. Instead, it interacts with an objective interface that returns
one scalar value for each queried point.

The objective interface can operate in either deterministic or stochastic
mode. In deterministic mode, the value returned to the optimizer corresponds
directly to the underlying clean objective. In stochastic mode, an
independent random perturbation is added to each objective query.

This design allows the same underlying optimization problem to be evaluated
under different fitness conditions without changing the benchmark function
itself. The stochasticity is introduced at the interface between the
generated optimizer and the objective rather than by modifying the
optimization algorithm directly.


%------------------------------------------------------------------------------
\subsection{Deterministic and Stochastic Evaluation}
\label{subsec:deterministic_stochastic}
%------------------------------------------------------------------------------

Let $f_{\mathrm{clean}}(x)$ denote the underlying benchmark objective
function for a candidate solution $x \in \mathbb{R}^{D}$. Let $x^*$ denote a
global optimum in the search space, and let

\begin{equation}
f^*
=
f_{\mathrm{clean}}(x^*)
\end{equation}

denote the corresponding optimal objective value. For deterministic
evaluation, the interface returns the clean objective value,

\begin{equation}
f_{\mathrm{noise}}(x)
=
f_{\mathrm{clean}}(x),
\end{equation}

which corresponds to the standard noiseless BBOB setting.

For stochastic evaluation, the returned value is instead given by

\begin{equation}
f_{\mathrm{noise}}(x)
=
f_{\mathrm{clean}}(x) + z,
\label{eq:stochastic_objective}
\end{equation}

where the perturbation is sampled independently for each objective query.
The standard deviation of the perturbation is scaled by the clean
optimality gap,

\begin{equation}
z
\sim
\mathcal{N}
\left(
0,
\left(
\sigma
\left|
f_{\mathrm{clean}}(x)-f^*
\right|
\right)^2
\right).
\label{eq:proportional_gaussian_noise}
\end{equation}

Equivalently, the dynamic standard deviation at a queried point is

\begin{equation}
\sigma_{\mathrm{dynamic}}(x)
=
\sigma
\left|
f_{\mathrm{clean}}(x)-f^*
\right|.
\label{eq:dynamic_noise_std}
\end{equation}

For $\sigma=0$, the stochastic term vanishes and the interface becomes
deterministic. For positive values of $\sigma$, the magnitude of the
perturbation changes with the current objective gap. The optimum and clean
objective are used internally by the evaluation wrapper to construct this
noise, but they are not exposed to the generated optimizer.

The use of the optimality gap rather than the raw objective value is
important because the same optimization problem can be shifted vertically
without changing its optimization structure. If noise were scaled directly
by the raw objective value, such a shift could change the magnitude of the
noise without changing the underlying optimization problem. Scaling by
$|f_{\mathrm{clean}}(x)-f^*|$ instead relates the stochastic perturbation to
the remaining optimization error.

The noise model used in this thesis therefore produces a heteroscedastic
objective: the magnitude of the perturbation depends on the objective gap.
This property is useful for studying optimization behavior because the
relative uncertainty changes as the optimizer approaches the optimum. An
algorithm can consequently experience different levels of stochastic
variation during different stages of its search.

For illustration, consider a point with an objective gap of approximately
$10^3$ and an illustrative noise parameter of $\sigma=0.1$. The dynamic
standard deviation is then

\begin{equation}
\sigma_{\mathrm{dynamic}}
=
0.1\cdot10^3
=
10^2.
\label{eq:f8_dynamic_std_boundary}
\end{equation}

If another point has an objective gap of approximately $10^1$, the
corresponding standard deviation becomes

\begin{equation}
\sigma_{\mathrm{dynamic}}
=
0.1\cdot10^1
=
10^0.
\label{eq:f8_dynamic_std_valley}
\end{equation}

These values are used only to illustrate the scaling behavior of the noise
model. The experimental noise settings are defined separately in
Chapter~\ref{chap:experiments}.

%------------------------------------------------------------------------------
\subsection{Illustration of the Noise Model}
\label{subsec:noise_model_illustration}
%------------------------------------------------------------------------------

The behavior of the noise scaling is illustrated using two BBOB functions
with different landscape structures. The first example is the Rosenbrock
function $f_8$, which contains a narrow curved valley. This structure makes
it useful for illustrating how the magnitude of the stochastic perturbation
changes between regions with different objective gaps
\cite{Hansen2010}.

The equations above show that a point with a larger objective gap receives a
larger perturbation, while a point with a smaller objective gap receives a
smaller perturbation. The scaling therefore depends on the objective gap and
not directly on the Euclidean distance of a point from the optimum.

\begin{figure}[htbp]
    \centering
    \includegraphics[width=\textwidth]{images/f8/figure_1d_noise_cross_section.png}

    \vspace{0.4cm}

    \includegraphics[width=\textwidth]{images/f8/figure_2d_noise_topology.png}

    \caption{Illustration of stochastic distortion on BBOB $f_8$
    (Rosenbrock). The one-dimensional cross-section (top) and two-dimensional
    contours (bottom) compare the clean objective with noisy evaluations.
    The optimality-gap model changes the noise magnitude according to the
    objective gap, resulting in smaller perturbations as the search approaches
    the optimum.}
    \label{fig:f8_noise_distortion}
\end{figure}

A second example is provided by BBOB $f_{21}$, the Gallagher 101-peaks
function. Unlike the narrow-valley structure of $f_8$, this function contains
many peaks and a highly irregular multimodal structure. It therefore provides
a different setting in which stochastic fitness can influence the ranking of
candidate points and consequently affect search decisions.

The noise model does not change the underlying clean function. Instead, each
query produces a noisy fitness value with a perturbation whose standard
deviation depends on the clean objective gap. Consequently, regions with
larger objective gaps can exhibit larger variation in their noisy fitness
values, while the magnitude of the perturbation decreases as the search
approaches the global optimum.

\begin{figure}[htbp]
    \centering
    \includegraphics[width=\textwidth]{images/f21/figure_1d_noise_cross_section.png}

    \vspace{0.4cm}

    \includegraphics[width=\textwidth]{images/f21/figure_2d_noise_topology.png}

    \caption{Illustration of stochastic distortion on BBOB $f_{21}$
    (Gallagher 101-peaks). The one-dimensional cross-section (top) and
    two-dimensional representation (bottom) illustrate the effect of
    stochastic objective evaluations on the multimodal function. The
    optimality-gap formulation changes the perturbation magnitude according
    to the clean objective gap while leaving the underlying benchmark
    function unchanged.}
    \label{fig:f21_noise_distortion}
\end{figure}

The use of two functions with different structures is intended to illustrate
that the stochastic evaluation mechanism is not tied to one particular
landscape shape. The noise is applied through the same objective wrapper,
while the resulting noisy fitness values depend on the clean objective
values of the function being evaluated.

%------------------------------------------------------------------------------
\subsection{Independent Objective Realizations}
\label{subsec:independent_noise_realizations}
%------------------------------------------------------------------------------

Every objective query receives a new stochastic realization. The same
stochastic objective model is therefore used for all queries, while the
random perturbation $z$ is independently sampled for each query. For the
same coordinate $x$, two evaluations can consequently return different
noisy fitness values,

\begin{equation}
f_{\mathrm{noise}}^{(i)}(x)
=
f_{\mathrm{clean}}(x) + z_i,
\qquad
z_i \overset{\mathrm{i.i.d.}}{\sim}
\mathcal{N}
\left(
0,
\left(
\sigma
\left|
f_{\mathrm{clean}}(x)-f^*
\right|
\right)^2
\right).
\label{eq:independent_noise_realizations}
\end{equation}

Thus, repeated queries at the same point can return different noisy fitness
values, even though the underlying clean objective and the stochastic
objective model remain unchanged.

The generated optimizer therefore has to operate on repeated noisy fitness
values rather than on a deterministic function with a modified scale. This
can affect algorithms that rely strongly on direct fitness comparisons,
ranking, selection, or changes in the returned objective values.

For example, if two candidate points have similar clean objective values, a
stochastic realization can temporarily reverse their ranking. A generated
optimizer that reacts immediately to such rankings may therefore make a
different search decision than it would under deterministic evaluation.
Repeated fitness values can similarly affect algorithms that estimate
progress from differences between successive objective values.

The stochastic environment thus changes the information available to the
optimizer without changing the underlying optimization problem. This
distinction is important for the experimental analysis because the objective
function itself remains fixed while the fitness values available during
search vary between objective queries.

%------------------------------------------------------------------------------
\subsection{Clean Diagnostic Evaluation}
\label{subsec:clean_diagnostic}
%------------------------------------------------------------------------------

The evaluation framework nevertheless retains the clean objective
internally. After an algorithm has completed its search, the returned
solution is evaluated using the clean objective rather than the particular
noisy fitness value returned during the final stochastic query.

For a returned solution $x_{\mathrm{best}}$, the clean optimality gap is
defined as

\begin{equation}
\Delta f(x_{\mathrm{best}})
=
\left|
f_{\mathrm{clean}}(x_{\mathrm{best}})
-
f^*
\right|.
\label{eq:clean_optimality_gap}
\end{equation}

This value provides a direct measure of the actual quality of the solution
returned by the generated optimizer. It is not exposed to the optimizer
during its execution.

The separation between the noisy fitness used during optimization and the
clean diagnostic objective is central to the methodology. The generated
algorithm has to make its search decisions using only the information
available through the stochastic objective interface. The synthesis
framework can then determine how good the resulting solution actually is
using the underlying clean objective.

This separation also prevents a favorable random fitness realization from
being interpreted as direct evidence that a candidate algorithm has produced
a better solution. The returned point is instead assessed according to its
clean objective value after the optimization run has finished.


%==============================================================================
\section{Candidate Selection and Evolutionary Update}
\label{sec:candidate_selection}
%==============================================================================

Once a generated candidate has been executed and evaluated, its result is
returned to the evolutionary controller. The controller uses the clean
diagnostic quality of the returned solution to compare the candidate with
the current incumbent. This is an important distinction from the fitness
values used internally by the generated optimizer: the candidate searches
using the configured objective interface, whereas the synthesis framework
compares the resulting algorithms using the clean optimality gap of their
returned solutions.

For a parent algorithm
$\mathcal{A}_{\mathrm{parent}}$ and a child algorithm
$\mathcal{A}_{\mathrm{child}}$, let

\begin{equation}
\Delta f_{\mathrm{parent}}
=
\Delta f
\left(
x_{\mathrm{parent}}
\right)
\end{equation}

and

\begin{equation}
\Delta f_{\mathrm{child}}
=
\Delta f
\left(
x_{\mathrm{child}}
\right)
\end{equation}

denote their respective clean optimality gaps.

Since the optimization problem is formulated as minimization, a child with a
smaller clean optimality gap represents an improvement over the parent. The
parent is therefore retained only when the child does not provide an
improvement according to the selection criterion.

%------------------------------------------------------------------------------
\subsection{Parent--Child Comparison}
\label{subsec:parent_child_comparison}
%------------------------------------------------------------------------------

The parent--child comparison represents the evolutionary update of the
synthesis process. The parent corresponds to the currently retained
algorithm, while the child is the newly generated algorithm produced from
the current synthesis context. The synthesis process uses one parent and one
offspring with elitist selection, corresponding to a $(1+1)$ evolutionary
structure.

Each generated child is executed once by the evaluation framework. During
this execution, the optimizer interacts with the configured objective
interface and receives the corresponding fitness values. Under stochastic
conditions, these values are subject to the configured noise model.

After the execution has finished, the returned solution is evaluated using
the clean objective. For a returned solution $x_{\mathrm{best}}$, the clean
optimality gap is defined as

\begin{equation}
\Delta f(x_{\mathrm{best}})
=
\left|
f_{\mathrm{clean}}(x_{\mathrm{best}})
-
f^*
\right|.
\label{eq:clean_optimality_gap_selection}
\end{equation}

The resulting error is converted into the fitness score used by LLaMEA,

\begin{equation}
S
=
-\Delta f,
\label{eq:synthesis_fitness_score}
\end{equation}

so that smaller clean optimality gaps correspond to higher fitness scores.
A successful candidate with an optimality gap of zero therefore receives the
maximum possible fitness score of zero.

The scored candidate is then passed to the LLaMEA synthesis engine. With one
parent, one offspring, and elitism enabled, the evolutionary update retains
the candidate with the better fitness score. Equivalently, in terms of the
clean optimality gap, the child replaces the incumbent only when its clean
error is strictly smaller than that of the currently retained parent:

\begin{equation}
\mathcal{A}_{\mathrm{inc}}^{(g+1)}
=
\begin{cases}
\mathcal{A}_{\mathrm{child}}^{(g)}
&
\text{if }
\Delta f_{\mathrm{child}}^{(g)}
<
\Delta f_{\mathrm{parent}}^{(g)},
\\[4pt]
\mathcal{A}_{\mathrm{parent}}^{(g)}
&
\text{otherwise}.
\end{cases}
\label{eq:clean_parent_child_selection}
\end{equation}

The parent is not re-executed before this comparison. Its previously
recorded clean diagnostic fitness remains the reference for elitist
selection.

This selection mechanism separates the stochastic search experienced by the
generated optimizer from the criterion used by the outer synthesis process.
The optimizer makes its search decisions using the fitness values returned
by the configured objective interface, whereas the resulting candidate
quality is assessed after execution using the clean objective. The clean
diagnostic result is then converted into the fitness score used by the
synthesis engine.

Although the clean objective is not available to the generated optimizer
during its execution, the resulting clean performance is included in the
feedback provided to the synthesis process. Thus, the stochastic objective
affects the candidate's search trajectory, while the clean diagnostic score
provides the performance signal used for evolutionary selection and
subsequent synthesis feedback.


%------------------------------------------------------------------------------
\subsection{Feedback and Evolutionary Context Update}
\label{subsec:feedback_update}
%------------------------------------------------------------------------------

After the parent--child comparison, the result of the candidate evaluation is
added to the synthesis context used in the following generation. This
feedback is not another selection criterion. The decision of whether the
child replaces the parent is still based on the clean objective error
described in the previous section. The role of the feedback is instead to
tell the LLM what happened during the evaluation so that this information can
be taken into account when generating the next candidate.

The feedback depends on how the generated algorithm behaved during
execution. A successfully executed candidate produces a result message,
whereas a candidate that cannot be executed produces diagnostic information
about the failure. Repeated execution failures can additionally trigger
meta-feedback that asks the LLM to change its search direction. For
stochastic objectives, diagnostic messages can also contain information
reminding the LLM that objective evaluations are stochastic and that
objective-call bookkeeping must be handled accordingly.

For a candidate that executes successfully, the evaluator reports its final
objective error. A typical result message is:

\begin{verbatim}
[RESULT]

The generated algorithm executed successfully.

Final objective error:

0.0012

Use this result together with the previous algorithm history when
designing the next candidate.
\end{verbatim}

When the candidate is evaluated under a stochastic objective, the feedback
can additionally make the stochastic setting explicit:

\begin{verbatim}
[RESULT]

The generated algorithm executed successfully on a stochastic objective.

Final objective error:

0.8420

The result indicates that the current search strategy may not have
handled the stochastic evaluations effectively.

Use the observed result and previous algorithm history to improve
the next candidate.
\end{verbatim}

The numerical error in these messages is obtained from the actual candidate
evaluation. The feedback does not tell the LLM which optimization algorithm
to use next. Instead, it gives the model the observed result and leaves the
next algorithmic decision to the synthesis process. This is important because
the feedback is intended to support the evolutionary generation process
without turning the evaluator into a hand-designed optimization strategy.

Candidates that fail during execution are handled differently. Rather than
providing only a poor score, the evaluator reports the reason for the
failure so that the next generated program can address it. For example, a
runtime error can be communicated as follows:

\begin{verbatim}
[RUNTIME ERROR]

The generated algorithm failed during execution.

Error:

ValueError: shapes (10, 2) and (3, 10) not aligned: 2 != 3

Relevant code:

  -> line   8:     arr = pop @ weights

Fix the cause of the failure in the next candidate.

Ensure that:

- array shapes and dimensions are valid;
- numerical operations are well-defined;
- all variables are initialized before use;
- the search stays within the provided bounds;
- every objective evaluation is counted;
- the total number of objective evaluations does not exceed the budget.

Problem context: BBOB-1, dim=3, bounds=[-5.0, 5.0].
\end{verbatim}

The message therefore gives the LLM information that is directly related to
the observed implementation problem. This can include invalid array
dimensions, undefined variables, numerical errors, violations of the search
bounds, or incorrect evaluation-budget handling. The intention is not to
explain the entire implementation to the model, but to provide enough
information to correct the failure in the next candidate.

For candidates evaluated under stochastic fitness, the runtime diagnostic
can be extended with the following context:

\begin{verbatim}
[NOISY PROBLEM CONTEXT]

The objective function is stochastic.

Ensure that optimization decisions are not based on invalid or
inconsistently stored objective values.

Keep any internal statistical estimates separate from the
objective value required by the optimizer interface.

Ensure that every call to problem(x) is counted against the
evaluation budget.
\end{verbatim}

This information is deliberately kept general. It does not require the LLM
to use a particular noise-handling technique, a fixed number of repeated
evaluations, or a predefined statistical estimator. Such decisions remain
part of the generated algorithm. The feedback only makes the relevant
properties of the execution environment clear when a diagnostic message
requires this context.

Execution time is treated separately from ordinary runtime errors. If a
candidate exceeds the allowed execution time, the evaluator reports:

\begin{verbatim}
[TIMEOUT]

The generated algorithm exceeded the execution time limit.

Reduce unnecessary computation and ensure that the implementation
can complete within the available execution time and evaluation
budget.

Problem context: BBOB-1, dim=3, bounds=[-5.0, 5.0].
\end{verbatim}

The stochastic problem context can again be attached when the candidate was
evaluated under noisy fitness. Other failures, such as syntax errors or
mathematical exceptions, are handled in the same general way: the evaluator
reports the observed problem and asks the LLM to correct the implementation
in the following candidate.

A separate case occurs when several consecutive candidates fail. Without
additional intervention, the LLM may continue making small changes to an
unsuccessful implementation. The current implementation therefore tracks
repeated failures and can add meta-feedback to the synthesis history:

\begin{verbatim}
[META-FEEDBACK]

The last several generated algorithms failed to execute correctly.

Try a substantially different search mechanism rather than making
only minor modifications to the previous approach.

Check the previous implementation for the source of the failure and
ensure that the new algorithm respects the objective interface,
search bounds, evaluation budget, and numerical constraints.
\end{verbatim}

The message deliberately does not recommend a particular optimization
algorithm. This keeps the synthesis process separate from the external
benchmark algorithms used later in the evaluation. The purpose of the
meta-feedback is simply to move the LLM away from repeatedly modifying the
same unsuccessful implementation.

Taken together, these messages provide the LLM with information about what
happened to the previous candidate while keeping the evolutionary selection
mechanism unchanged. A successful candidate contributes its observed
performance, a failed candidate contributes information about the execution
problem, and repeated failures can trigger a request for a substantially
different search mechanism. This information is accumulated in the
synthesis history represented by $\mathcal{H}_{1:g-1}$ in
Equation~\ref{eq:synthesis_context}. At generation $g$, the relevant history
is provided to the LLM together with the current task prompt. The feedback
therefore connects candidate evaluation with the generation of the next
candidate, while the parent--child selection rule remains based on the
underlying scalar objective error.


%------------------------------------------------------------------------------
\subsection{Prompt Configurations}
\label{subsec:prompt_configurations}
%------------------------------------------------------------------------------

The prompt given to the LLM is assembled from three components:
\texttt{TASK}, \texttt{EXAMPLE}, and \texttt{FORMAT}. This separation was
introduced so that the information about the optimization environment can be
changed without changing the required program interface or the general code
structure expected from the model. The \texttt{EXAMPLE} and
\texttt{FORMAT} components are kept identical for all synthesis conditions,
while the \texttt{TASK} component changes according to the experimental
condition.

The complete prompt can therefore be represented as

\begin{equation}
\mathrm{Prompt}
=
\mathrm{TASK}
+
\mathrm{EXAMPLE}
+
\mathrm{FORMAT}.
\label{eq:prompt_composition}
\end{equation}

The task itself is assembled from a common layout, an environment mode, and
a strategy scaffold:

\begin{equation}
\mathrm{TASK}
=
\mathrm{Layout}
+
\mathrm{Environment\ Mode}
+
\mathrm{Strategy\ Scaffold}.
\label{eq:prompt_task_composition}
\end{equation}

This structure gives two independent factors that can be varied during
synthesis. The first factor controls what the LLM is told about the
evaluation environment. The second controls the additional guidance given
about the search strategy or implementation. There are three environment
modes, \textsc{Clean}, \textsc{Implicit}, and \textsc{Noisy}, and four
strategy scaffolds, \textsc{Baseline}, \textsc{Vectorization},
\textsc{Guided}, and \textsc{Thinking}. Their combination results in twelve
different prompt conditions.

Figure~\ref{fig:prompt_architecture} provides an overview of this modular prompt architecture, showing how the invariant skeleton (\texttt{FORMAT}, \texttt{EXAMPLE}, and problem \texttt{Layout}) is composed with the two-factor design.

\begin{figure}[htbp]
\centering
\resizebox{\textwidth}{!}{%
\begin{tikzpicture}[
    font=\footnotesize,
    >=Stealth,
    % -------------------------------------------------------------------------
    % Node Styling
    % -------------------------------------------------------------------------
    base/.style={
        rectangle, 
        rounded corners=3pt, 
        thick, 
        align=center, 
        inner sep=3.5pt
    },
    root/.style={
        base,
        draw=blue!75!black,
        fill=blue!10,
        text width=9.8cm,
        minimum height=0.7cm,
        font=\small\bfseries
    },
    fixedbox/.style={
        base,
        draw=blue!75!black,
        fill=blue!8,
        text width=3.3cm,
        minimum height=0.65cm,
        font=\scriptsize\bfseries
    },
    fixeditem/.style={
        base,
        rounded corners=2pt,
        draw=blue!50!black,
        fill=blue!2,
        align=left,
        text width=3.3cm,
        inner sep=3.5pt,
        font=\tiny
    },
    variablebox/.style={
        base,
        draw=orange!75!black,
        fill=orange!8,
        text width=3.3cm,
        minimum height=0.65cm,
        font=\scriptsize\bfseries
    },
    modeparent/.style={
        base,
        rounded corners=3pt,
        text width=2.9cm,
        minimum height=0.65cm,
        font=\scriptsize\bfseries
    },
    subenv/.style={
        base,
        rounded corners=2pt,
        text width=1.35cm,
        minimum height=0.55cm,
        font=\scriptsize
    },
    stratbanner/.style={
        base,
        draw=purple!65!black,
        fill=purple!8,
        text width=7.4cm,
        minimum height=0.55cm,
        font=\scriptsize\bfseries
    },
    stratbox/.style={
        base,
        rounded corners=2pt,
        draw=purple!60!black,
        fill=purple!5,
        text width=1.7cm,
        minimum height=0.55cm,
        inner sep=2pt,
        font=\scriptsize
    },
    treearrow/.style={
        ->,
        >=Stealth,
        thick,
        draw=black!70
    }
]

    % =========================================================================
    % 1. ROOT
    % =========================================================================
    \node[root] (root) {
        Prompt Specification\\[1pt]
    };

    % =========================================================================
    % 2. TOP LEVEL (INVARIANT vs. FACTORIAL)
    % =========================================================================
    \node[
        fixedbox, 
        below=1.1cm of root.south,
        xshift=-3.9cm
    ] (fixed) {
        Invariant Prompts\\[1pt]
    };

    \node[
        variablebox, 
        below=1.1cm of root.south,
        xshift=1.9cm
    ] (variable) {
        Variant Prompts
    };

    \draw[treearrow] (root.south) -- ++(0,-0.4) -| (fixed.north);
    \draw[treearrow] (root.south) -- ++(0,-0.4) -| (variable.north);

    % =========================================================================
    % 3. LEFT COLUMN: BUS-STYLE ROUTED STACK
    % =========================================================================
    \node[fixeditem, below=0.85cm of fixed] (layout) {
        \textbf{Task Layout}\\
        Target identifier, dimension $D$, bounds $[-5, 5]^D$, evaluation budget $B_{\text{synth}}$, and tags.
    };

    \node[fixeditem, below=0.85cm of layout] (example) {
        \textbf{\texttt{EXAMPLE}}\\
        Common algorithm class, \texttt{\_\_init\_\_}, \texttt{\_\_call\_\_}, bound extraction, budget tracker.
    };

    \node[fixeditem, below=0.85cm of example] (format) {
        \textbf{\texttt{FORMAT}}\\
        Uniform response schema, single block, return tuple $(\mathbf{x}^*, y_{\text{raw}})$, whitelist NumPy.
    };

    % Clean, dedicated left bus line so arrows don't shoot straight through nodes
    \draw[treearrow, rounded corners=2pt] (fixed.west) -- ++(-0.35,0) |- (layout.west);
    \draw[treearrow, rounded corners=2pt] (fixed.west) -- ++(-0.35,0) |- (example.west);
    \draw[treearrow, rounded corners=2pt] (fixed.west) -- ++(-0.35,0) |- (format.west);

    % =========================================================================
    % 4. RIGHT COLUMN: IMPLICIT & EXPLICIT (CO-EQUAL LEVEL)
    % =========================================================================
    \node[
        modeparent,
        draw=yellow!65!black,
        fill=yellow!12,
        below=0.8cm of variable.south,
        xshift=-1.9cm
    ] (implicit) {
        \textbf{Implicit Mode}\\[1pt]
        \tiny Incomplete context; evaluation variability undisclosed
    };

    \node[
        modeparent,
        draw=orange!75!black,
        fill=orange!12,
        below=0.8cm of variable.south,
        xshift=1.9cm
    ] (explicit) {
        \textbf{Explicit Mode}\\[1pt]
        \tiny Evaluation behavior is explicitly characterized
    };

    \draw[treearrow] (variable.south) -- ++(0,-0.3) -| (implicit.north);
    \draw[treearrow] (variable.south) -- ++(0,-0.3) -| (explicit.north);

    % =========================================================================
    % 5. EXPLICIT SUB-BRANCHES: CLEAN & NOISY
    % =========================================================================
    \node[
        subenv,
        draw=green!60!black,
        fill=green!8,
        below=0.75cm of explicit.south,
        xshift=-0.85cm
    ] (clean) {
        \textbf{Clean}
    };

    \node[
        subenv,
        draw=red!60!black,
        fill=red!8,
        below=0.75cm of explicit.south,
        xshift=0.85cm
    ] (noisy) {
        \textbf{Noisy}
    };

    \draw[treearrow] (explicit.south) -- ++(0,-0.25) -| (clean.north);
    \draw[treearrow] (explicit.south) -- ++(0,-0.25) -| (noisy.north);

    % =========================================================================
    % 6. STRATEGY SCAFFOLD BANNER (CLEAR VERTICAL GAP BELOW CLEAN/NOISY)
    % =========================================================================
    \node[
        stratbanner,
        below=0.9cm of clean.south -| variable.south
    ] (strat_title) {
        \textbf{Strategy Scaffold}
    };

    % Direct vertical drops into the banner
    \draw[treearrow] (implicit.south) -- (implicit.south |- strat_title.north);
    \draw[treearrow] (clean.south)    -- (clean.south |- strat_title.north);
    \draw[treearrow] (noisy.south)    -- (noisy.south |- strat_title.north);

    % =========================================================================
    % 7. FOUR STRATEGY BLOCKS (GENEROUS BOTTOM MARGIN)
    % =========================================================================
    \node[stratbox, below=0.6cm of strat_title.south west, anchor=north west, xshift=0.08cm] (s_base) {
        \textbf{Baseline}
    };

    \node[stratbox, right=0.18cm of s_base] (s_think) {
        \textbf{Thinking}
    };

    \node[stratbox, right=0.18cm of s_think] (s_vec) {
        \textbf{Vectorization}
    };

    \node[stratbox, right=0.18cm of s_vec] (s_guided) {
        \textbf{Guided}
    };

    \draw[treearrow] (strat_title.south -| s_base.north)   -- (s_base.north);
    \draw[treearrow] (strat_title.south -| s_think.north)  -- (s_think.north);
    \draw[treearrow] (strat_title.south -| s_vec.north)    -- (s_vec.north);
    \draw[treearrow] (strat_title.south -| s_guided.north) -- (s_guided.north);

\end{tikzpicture}
}
\caption{Modular prompt architecture and factorial synthesis design. The structural components (\texttt{FORMAT}, \texttt{EXAMPLE}, and problem \texttt{Layout}) remain fixed across all synthesis runs to guarantee executable syntax and uniform evaluation. The task prompt is varied across two independent factors: three environment information modes (\textsc{Clean}, \textsc{Implicit}, and \textsc{Noisy}) and four search strategy scaffolds (\textsc{Baseline}, \textsc{Vectorization}, \textsc{Guided}, and \textsc{Thinking}), yielding twelve distinct experimental prompt conditions.}
\label{fig:prompt_architecture}
\end{figure}

The common \texttt{FORMAT} component defines how the LLM must return the
generated algorithm. Its complete content is:

\begin{verbatim}
Respond with EXACTLY the following format — no extra code blocks:

Feedback: <your reasoning and description of the algorithm>

Code:

```python

<your complete class and any required imports>

````

STRICT Rules — violating any rule will cause execution failure:

* There must be exactly ONE `python ... ` block in your response.

* The class MUST be named exactly one word (e.g., class MyOptimizer:).

* **init**(self) MUST take NO extra arguments beyond self.

* **init**(self) MUST have a non-empty body (use pass if nothing to initialize).

* The class MUST have a **call**(self, problem, budget) method.

* **call** MUST return (best_x, float(best_y)) — a tuple of the best
  search coordinates array and best scalar float value.

* Do NOT import or call scipy.optimize (e.g. scipy.optimize.minimize,
  differential_evolution, etc.) — pre-built solver wrappers are strictly
  banned. Write your search algorithm logic from scratch using NumPy.

* Every variable you use MUST be defined before use. Never reference
  undefined names.

* Do NOT store problem or budget in **init** — they are provided to
  **call** directly.

* Do NOT include if **name** == '**main**': blocks.
\end{verbatim}

The output requirements are common to every generated algorithm. In
particular, the return value must always contain \texttt{best\_x} and \texttt{best\_y}. These are part of the common optimizer interface and are
therefore not repeated in the environment-specific prompt components. This
keeps the interface identical across all synthesis conditions.

The format also requires the optimization logic to be written by the LLM
rather than delegated to a pre-built optimization solver. In particular,
\texttt{scipy.optimize} is prohibited. The generated implementation must
instead provide the search procedure directly and must respect the objective
evaluation budget.

The second invariant component is \texttt{EXAMPLE}. It provides a common
code skeleton showing how the generated optimizer is instantiated and
called:

\begin{verbatim}
Your algorithm will be instantiated and called as follows:

```
optimizer = AlgorithmName()

best_x, best_y = optimizer(problem, budget)
```

You MUST use the following class skeleton — fill in your algorithm
logic in the marked section only.

Do NOT change the class structure, method signatures, or return statement:

```
import numpy as np

class AlgorithmName:

    def __init__(self):

        pass  # Add initialization state here if your algorithm needs it

    def __call__(self, problem, budget):

        lb = np.asarray(getattr(problem, 'lower_bound', -5.0), dtype=float)

        ub = np.asarray(getattr(problem, 'upper_bound', 5.0), dtype=float)

        dim = int(getattr(problem, 'dim', len(lb) if hasattr(lb, '__len__') else 3))

        # Always start with a random initial point using vectorization

        best_x = np.random.uniform(lb, ub, size=dim)

        best_y = float(problem(best_x))

        evaluations = 1

        # --- YOUR ALGORITHM LOGIC BELOW ---

        # Use `evaluations` to track calls. Stop when evaluations >= budget.

        # BUDGET TRACKING: every call to problem(x) counts — including calls inside
        # min()/max(key=...) lambdas, list comprehensions, and helper functions.

        # Always increment `evaluations` immediately after every problem(x) call.

        # Leverage NumPy matrix vectorization across populations of shape (pop_size, dim).

        # Example batch mutation: candidates = np.clip(pop + np.random.normal(0, 0.1, size=(pop_size, dim)), lb, ub)

        # Evaluate each candidate: for x in candidates: y = float(problem(x)); evaluations += 1

        # Compare candidates and update best_x, best_y when improvement found.

        # NOTE: For noisy problems, a single raw comparison may be unreliable.

        # --- YOUR ALGORITHM LOGIC ABOVE ---

        return best_x, float(best_y)  # MUST return tuple: (best_x_ndarray, best_y_float)
```

\end{verbatim}

The example fixes the interface, initialization, access to the problem
bounds, and the evaluation counter. It also shows how population-based
operations can be implemented using NumPy. It does not, however, require
the generated algorithm to use a particular optimization strategy. The
actual search mechanism is left to the LLM and can therefore differ between
generated candidates.

The variable part of the prompt follows a common task layout:

\begin{verbatim}
You are designing a continuous black-box optimization algorithm.

The algorithm will be evaluated on the specified optimization problem
using only calls to problem(x).

Problem information:

* BBOB function ID: {{ problem_id }}

* Dimension: {{ dimension }}

* Lower bound: {{ lower_bound }}

* Upper bound: {{ upper_bound }}

* Evaluation budget: {{ budget }}

{{ mode_prompt }}

{{ strategy_prompt }}

Design an effective optimization algorithm for this setting.

The algorithm must respect the provided search bounds and evaluation
budget.
\end{verbatim}

The problem identifier, dimension, bounds, and budget are filled in by the
synthesis system for each run. The two remaining placeholders contain the
selected environment mode and strategy scaffold.

The first environment condition is \textsc{Clean}. It explicitly tells the
LLM that the objective is deterministic:

\begin{verbatim}
Environment:

The objective function is deterministic. Repeated evaluations at
the same point return consistent objective values.

You do not need to account for stochastic evaluation noise when
making optimization decisions.
\end{verbatim}

The second condition is \textsc{Implicit}. Here, the LLM is not explicitly
told that the objective is noisy or deterministic. Instead, the prompt states
that repeated evaluations may differ and that no further information about
the source or magnitude of this variation is available:

\begin{verbatim}
Environment:

The objective function may return different values when evaluated at
the same point.

No further information about the source or magnitude of this
variation is available.

Design the algorithm using only the information available through
objective evaluations.
\end{verbatim}

The third condition is \textsc{Noisy}. In this case, stochasticity is stated
directly and the LLM is told why a single observation may not be sufficient
for a reliable comparison:

\begin{verbatim}
Environment:

The objective function is stochastic. Evaluating the same point
multiple times may produce different objective values because of
random variation.

Therefore, a single objective evaluation may be an unreliable basis
for deciding which candidate is better.

Design the algorithm so that its optimization decisions account for
this stochasticity.
\end{verbatim}

The three environment descriptions therefore provide different amounts of
information about the objective evaluation process. The underlying
optimization problem and execution interface remain the same.

The second prompt factor is the strategy scaffold. The
\textsc{Baseline} condition does not add any strategy-specific guidance.
This provides the least constrained version of the task and leaves the
choice of search mechanism entirely to the LLM.

The \textsc{Vectorization} condition adds implementation-oriented guidance:

\begin{verbatim}
Strategy guidance:

Use population-based representations where they are appropriate
for the optimization problem.

Represent multiple candidate solutions together and use NumPy array
operations for candidate generation, transformation, and population
updates where practical.

The implementation should remain within the available evaluation
budget.
\end{verbatim}

The intention here is to encourage efficient array-based operations when
they are useful, without fixing a population size or a particular update
rule.

The \textsc{Guided} condition provides more general search guidance:

\begin{verbatim}
Strategy guidance:

Design the search around a clear balance between exploration and
exploitation.

Use information from previously evaluated candidate solutions to
guide subsequent search steps.

Consider maintaining and updating multiple candidate solutions when
this supports the search strategy.

Adapt the search behavior as the optimization progresses rather than
using a fixed search step throughout the entire budget.

Use the available evaluation budget deliberately, allocating
evaluations between discovering promising regions and refining
promising solutions.

The algorithm should use objective evaluations as the primary source
of information and should not assume access to gradients or internal
properties of the objective function.
\end{verbatim}

Unlike a prescription for a particular optimizer, this scaffold describes
general properties that can be used when designing a search procedure. The
LLM remains responsible for deciding how these ideas should be implemented.

The final scaffold is \textsc{Thinking}. Instead of supplying additional
optimization rules, it asks the LLM to consider the main design decisions
before generating the implementation:

\begin{verbatim}
Strategy guidance:

Before writing the code, briefly reason about the main search
mechanism, how candidate solutions will be generated and selected,
and how the evaluation budget will be allocated.

The implementation should directly reflect this reasoning.

Keep the reasoning focused on decisions that affect the algorithm
rather than explaining general optimization concepts.
\end{verbatim}

This condition therefore changes the requested generation process rather than
specifying a particular optimization method.

The two prompt factors are combined to form the complete set of synthesis
conditions:

\begin{equation}
3\ \text{environment modes}
\times
4\ \text{strategy scaffolds}
=
12\ \text{prompt conditions}.
\label{eq:prompt_factorial_design}
\end{equation}

The resulting configuration matrix is illustrated in
Figure~\ref{fig:prompt_architecture} and summarized in
Table~\ref{tab:prompt_configuration_matrix}.

\begin{table}[htbp]
\centering
\caption{Prompt configuration matrix used for algorithm synthesis.}
\label{tab:prompt_configuration_matrix}
\begin{tabular}{lcccc}
\hline
Environment & Baseline & Vectorization & Guided & Thinking \\
\hline
Clean    & \checkmark & \checkmark & \checkmark & \checkmark \\
Implicit & \checkmark & \checkmark & \checkmark & \checkmark \\
Noisy    & \checkmark & \checkmark & \checkmark & \checkmark \\
\hline
\end{tabular}
\end{table}

For example, the \textsc{Clean--Baseline} condition combines the
deterministic environment description with no additional strategy guidance,
whereas \textsc{Noisy--Thinking} combines the explicit stochastic
environment description with the reasoning scaffold. The same
\texttt{FORMAT} and \texttt{EXAMPLE} components are used in both cases.

This separation is useful for interpreting the generated algorithms because
the two factors can be changed independently. Changing the strategy
scaffold while keeping the environment mode fixed changes the amount or type
of guidance given to the LLM. Changing the environment mode while keeping
the strategy scaffold fixed changes the information available to the model
about the objective evaluation process.

The prompt configurations do not prescribe external benchmark algorithms.
Methods such as CMA-ES, Differential Evolution, and particle swarm
optimization are used separately as reference algorithms during independent
evaluation and are not presented to the LLM as required solutions. Likewise,
the prompt does not impose a fixed number of repeated evaluations, a fixed
statistical estimator, or a fixed fraction of the evaluation budget for
handling stochastic fitness. These choices remain part of the generated
algorithm.

The complete prompt design therefore keeps the program interface,
output requirements, and common code example fixed while systematically
varying the information about the optimization environment and the level of
additional strategy guidance. The twelve resulting combinations form the
prompt conditions used by the synthesis procedure.


%------------------------------------------------------------------------------
\subsection{Separation of Noisy Search and Clean Selection}
\label{subsec:noisy_search_clean_selection}
%------------------------------------------------------------------------------

The synthesis process separates the stochastic fitness information available
to a generated optimizer from the clean objective used by the outer synthesis
controller. This separation is important because the generated optimizer may
make search decisions based on noisy fitness values, while the synthesis
controller must determine whether the resulting candidate should replace the
current incumbent.

During candidate execution, the optimizer interacts only with the objective
interface provided to it. Under stochastic conditions, repeated evaluations
of the same solution can therefore produce different fitness values. These
values influence the optimizer's internal search decisions and can affect
which solution it returns at the end of its evaluation budget.

After execution, the returned solution is evaluated separately using the
clean objective. The resulting clean optimality gap
$\Delta f(x_{\mathrm{best}})$ is used to assess the quality of the generated
candidate. The parent is represented by its previously recorded clean score
and is not re-executed before the comparison.

This creates a clear separation between the two roles of fitness information.
Noisy fitness determines the optimization environment experienced by the
generated algorithm, whereas clean fitness provides the common diagnostic
criterion used by the synthesis controller for parent--child selection.
Consequently, stochastic fitness can influence the candidate's search
trajectory and the solution it returns without directly determining whether
that candidate replaces the current incumbent.

%==============================================================================
\section{Independent Benchmark Evaluation}
\label{sec:independent_benchmark}
%==============================================================================

The synthesis procedure and the subsequent benchmark evaluation have
different roles and are therefore treated separately.

During synthesis, the objective is to generate and refine executable
optimization algorithms. Evaluation results are used by the evolutionary
process to determine whether a newly generated candidate should become the
current incumbent. The synthesis environment therefore directly influences
the algorithm-generation process.

Once the synthesis process has produced the retained algorithms, they are
evaluated independently using a common benchmark protocol. At this stage,
the synthesized algorithms are fixed and no further evolutionary modification
takes place. The purpose of the independent evaluation is to measure their
performance under predefined optimization conditions and to provide a common
basis for comparison with established optimization methods.

The separation between synthesis and benchmark evaluation also enables
cross-environment comparisons. An algorithm synthesized under deterministic
fitness can, for example, be evaluated under stochastic fitness, while an
algorithm synthesized under stochastic fitness can be evaluated under
deterministic fitness. This makes it possible to examine whether properties
acquired during synthesis remain effective when the evaluation environment
changes.

%------------------------------------------------------------------------------
\subsection{Benchmark Evaluation Protocol}
\label{subsec:benchmark_protocol}
%------------------------------------------------------------------------------

The independent benchmark evaluation applies the same general evaluation
procedure to the synthesized algorithms and the selected baseline
optimizers. Each algorithm is executed independently for each benchmark
condition, using predefined random seeds and an objective-function evaluation
budget. The resulting runs provide performance data that can be used to
analyze convergence, final solution quality, and robustness.

The evaluation budget is defined in terms of objective-function evaluations,
providing a common resource measure for algorithms with different internal
search procedures. Under stochastic conditions, the optimizer interacts with
the noisy objective through the same objective interface used during
synthesis, while final solution quality is assessed using the corresponding
clean objective.

Importantly, the benchmark evaluation is independent of the synthesis
process. The benchmark results do not modify the synthesized algorithms or
feed back into the evolutionary generation procedure. This allows the
performance of algorithms produced under different synthesis conditions to
be compared under the same evaluation protocol.

The specific benchmark functions, dimensions, noise levels, evaluation
budgets, repetitions, synthesized model variants, baseline algorithms,
evaluation targets, and statistical procedures are defined in
Chapter~\ref{chap:experiments}. Keeping these experimental parameters
separate from the methodology makes it possible to distinguish the general
operation of the synthesis and evaluation framework from the particular
conditions used in the experiments.